% =====================================================================
% Rozdział 6 – Podsumowanie (~2 str.)
% Nie streszczenie i nie parafraza celu: konkretne wyniki, co się nie
% udało, kierunki rozwoju.
% =====================================================================
\chapter{Podsumowanie}
\label{ch:podsumowanie}

% ---------------------------------------------------------------------
\section{Wnioski}
\label{sec:wnioski}
% - Konkretne wyniki przy każdym celu szczegółowym z \ref{sec:cel-zakres}.
% - Najważniejsze liczby: ok. 9 FPS na kamerę, 171 ms od zdjęcia do UDP,
%   dokładność azymutu, skuteczność detekcji.
% - Co się nie udało / czego nie zrealizowano i dlaczego.
% Stan: ⏳ (na koniec)


% ---------------------------------------------------------------------
\section{Kierunki dalszego rozwoju}
\label{sec:rozwoj}
% - Profil czasu warstw sieci (przygotowany, niewykonany) – dlaczego sieć
%   wolniejsza niż w benchmarkach Espressif.
% - Runtime esp-tflite-micro zamiast TFLite Micro z Edge Impulse.
% - Kod DSP poza -O0.
% - Przygotowanie klatki na drugim rdzeniu (równoległość z inferencją);
%   dwie instancje modelu na dwóch rdzeniach.
% - Wejście modelu bezpośrednio z uint8 (krok 6).
% - Więcej klas obiektów, szacowanie odległości, śledzenie obiektów
%   w czasie.
% Stan: ✅

